%------------------------------------------------------------
% 3.7 Comparación F_c vs F_g entre protón y electrón (átomo de hidrógeno)
%------------------------------------------------------------
\subsection*{3.7 \textcolor{Cinnabar}{\textbf{Comparación entre la fuerza eléctrica y la gravitacional}}}

Consideremos un protón y un electrón separados por una distancia típica del átomo de hidrógeno
\(r = 5.29\times 10^{-11}\ \text{m}\) (radio de Bohr). Tomamos
\(q_p=+e\), \(q_e=-e\) y masas \(m_p, m_e\).

\paragraph{Constantes:}
\(
k=\dfrac{1}{4\pi\varepsilon_0}=8.99\times 10^9\ \text{N m}^2\!/\text{C}^2,\;
e=1.602\times 10^{-19}\ \text{C},\;
G=6.67\times 10^{-11}\ \text{SI},\;
m_p=1.673\times 10^{-27}\ \text{kg},\;
m_e=9.109\times 10^{-31}\ \text{kg}.
\)

\paragraph{Fuerza de Coulomb:}
\[
\vec F_c = k\,\frac{q_p q_e}{r^2}\,\hat r
\quad\Rightarrow\quad
|\!F_c\!| = k\,\frac{e^2}{r^2}
          \approx \boxed{8.244e-08\ \text{N}}.
\]

\paragraph{Fuerza gravitacional (Newton):}
\[
\vec F_g = G\,\frac{m_p m_e}{r^2}\,\hat r
\quad\Rightarrow\quad
|\!F_g\!| = G\,\frac{m_p m_e}{r^2}
          \approx \boxed{3.634e-47\ \text{N}}.
\]

\paragraph{Cociente de magnitudes:}
\[
\boxed{\frac{F_c}{F_g} \;=\; \frac{k\,e^2}{G\,m_p m_e}
\;\approx\; 2.27e+39}
\;\;\;\approx\;\; 2.3\times 10^{39}.
\]

\paragraph{Conclusión:}
La fuerza eléctrica es del orden de \(10^{39}\) veces \emph{más intensa} que la fuerza
gravitacional entre un protón y un electrón a escala atómica; por ello, a nivel subatómico,
la interacción gravitacional es prácticamente \textbf{despreciable} frente a la de Coulomb.
